\begin{table}[htbp]
  \centering
  \caption{Espacios funcionales continuos y discretización en FEniCS.}
  \label{tab:functional_spaces}
  \begin{tabular}{p{3.2cm} c c p{4.2cm} p{2.8cm} c}
    \hline\hline
    \textbf{Variable Física} & \textbf{Símbolo} & \textbf{Espacio Continuo} & \textbf{Espacio Discreto (FEniCS)} & \textbf{Base Polinomial} & \textbf{Dimensión / DoFs} \\ \hline
    Potencial Eléctrico & $u_j$ & $H^1(\Omega)$ & FunctionSpace(mesh, 'P', 1) & Lagrange lineal continuo & 1089 nodos \\ \hline
    Conductividad Cuadrantes & $\sigma$ & $L^\infty(\Omega)$ & FunctionSpace(mesh, 'DG', 0) & Constante por trozos & 4 subdominios (L=4) \\ \hline
    Densidad Corriente Normal & $g_j = -I/|\mathcal{E}_j|$ & $H^{-1/2}(\partial\Omega)$ & Measure('ds', subdomain_data) & Constante por electrodo & 3 canales activos \\ \hline
    Flujo de Corriente & $\mathbf{J} = -\sigma \nabla u$ & $[L^2(\Omega)]^2$ & VectorFunctionSpace(mesh, 'P', 1) & Vectorial proyectado & 2178 componentes \\ \hline
    Malla de Dominio & $\Omega = [-1,1]^2$ & $\mathbb{R}^2$ & RectangleMesh & Triángulos estructurados & 2048 celdas \\ \hline
    \hline\hline
  \end{tabular}
\end{table}
